\section{The routing density and repeated irreducible copies}\label{sec:network}
A sweep followed by its reflection has a recursive description: an outer layer of fair switches, two independent half-sized networks, and another outer layer. The architecture is the Bene\v{s} network~\cite{Benes1964}. Its two-child probabilistic recursion was studied in~\cite[Section 4, Proposition 1]{GelmanTaShma2014}; we derive below the partial-permutation density identity needed here. The goal is to express a partial-permutation density in terms of the cycles created by its routing constraints. We then prove a truncation lemma that converts a density bound into contraction on repeated copies of an irreducible representation.

\subsection{Alternating colors and multiplicity}
Fix an ordered list of $k$ input positions and a candidate list of $k$ distinct output positions. At one node, make a graph on its tracked paths, joining two when they use the same outer input switch or the same outer output switch. Its degree is at most two. Input and output edges alternate; parallel edges form a cycle. If $s$ paths are present, $a$ is the number of edges, and $C$ is the number of cycles, the number of proper assignments of paths to the two children is
\[
 2^{s-a+C}.
\]
Indeed a path component has one more vertex than edge, whereas a cycle
has equally many. Each component admits two child assignments, obtained
by exchanging the colors. Thus each cycle supplies one factor of two
beyond $2^{s-a}$, as illustrated in Figure~\ref{fig:network-node}.

\begin{figure}[ht]
\centering
\begin{tikzpicture}[x=1cm,y=1cm,>=Stealth,font=\small]
\node[draw,rounded corners,minimum width=2.5cm,minimum height=1.2cm] (a) at (0,1.1) {child network $0$};
\node[draw,rounded corners,minimum width=2.5cm,minimum height=1.2cm] (b) at (0,-1.1) {child network $1$};
\foreach \y in {1.5,.7,-.7,-1.5}{
 \draw[->] (-4,\y)--(-3,\y);
 \draw[->] (3,\y)--(4,\y);
}
\draw[thick] (-3,1.5)--(-3,.7);
\draw[thick] (-3,-.7)--(-3,-1.5);
\draw[thick] (3,1.5)--(3,.7);
\draw[thick] (3,-.7)--(3,-1.5);
\draw[blue] (-3,1.5)--(-1.25,1.3);
\draw[red] (-3,.7)--(-1.25,-.9);
\draw[blue] (-3,-.7)--(-1.25,.9);
\draw[red] (-3,-1.5)--(-1.25,-1.3);
\draw[blue] (1.25,1.3)--(3,1.5);
\draw[blue] (1.25,.9)--(3,-.7);
\draw[red] (1.25,-.9)--(3,.7);
\draw[red] (1.25,-1.3)--(3,-1.5);
\node[align=center] at (-3,2.3) {input pair\\switches};
\node[align=center] at (3,2.3) {output pair\\switches};

\begin{scope}[
 tracked/.style={circle,draw,minimum size=5mm,inner sep=0pt},
 child zero/.style={tracked,fill=blue!10},
 child one/.style={tracked,fill=red!10},
 input constraint/.style={thick},
 output constraint/.style={thick,dashed}
]
\node at (0,-2.2) {Example components: $s$ tracked paths, $a$ edges};
\node[child zero] (p0) at (-4,-3.25) {$0$};
\node[child one]  (p1) at (-2.8,-3.25) {$1$};
\node[child zero] (p2) at (-1.6,-3.25) {$0$};
\node[child one]  (p3) at (-.4,-3.25) {$1$};
\draw[input constraint] (p0)--node[above,font=\scriptsize] {input} (p1);
\draw[output constraint] (p1)--node[above,font=\scriptsize] {output} (p2);
\draw[input constraint] (p2)--node[above,font=\scriptsize] {input} (p3);
\node[child zero] (c0) at (1.4,-3.25) {$0$};
\node[child one]  (c1) at (3.4,-3.25) {$1$};
\draw[input constraint] (c0) to[bend left=35]
 node[above,font=\scriptsize] {input} (c1);
\draw[output constraint] (c0) to[bend right=35]
 node[below,font=\scriptsize] {output} (c1);
\node[align=center] at (-2.2,-4.7)
 {path: $s=4$, $a=3$\\$2=2^{s-a}$ child assignments};
\node[align=center] at (2.4,-4.7)
 {cycle: $s=2$, $a=2$\\$2=2^{s-a}\cdot 2$ child assignments};
\node[align=center] at (0,-5.75)
 {Each edge forces opposite child labels.\\On each component, swapping $0$ and $1$ gives the second assignment.};
\end{scope}
\end{tikzpicture}

\caption{A recursive network node (top) and two components of a routing-constraint graph (bottom). Each graph vertex represents a tracked path, and each input or output edge forces its endpoints into opposite children. Both displayed components have two child assignments, but the cycle contributes an extra factor $2$ relative to $2^{s-a}$. The four paths in the network panel are schematic; each child contains half the positions.}
\label{fig:network-node}
\end{figure}

Each assignment prescribes $2s-a$ independent outer coins, hence costs $2^{-(2s-a)}$. Their combined factor is $2^{-s+C}$. Recursing through the nodes and choosing proper colorings uniformly gives a density relative to the uniform injection with excess exponent
\begin{equation}\label{net:exponent}
 X=\sum_{j=1}^d C_j+\log_2\frac{(N)_k}{N^k},
\end{equation}
where $C_j$ totals tracked cycles at height $j$. We distinguish the auxiliary coloring law from the actual switch law. For fixed input and output injections $x,y$, let $\omega$ be a legal complete tree of child colorings. At a node $v$, write $s_v,e_v,C_v$ for its path, edge and cycle counts, and put $a_v=2^{s_v-e_v+C_v}$. The auxiliary weight and actual routing-event weight are respectively
\[
 q_{xy}(\omega)=\prod_v a_v^{-1},\qquad
 w_{xy}(\omega)=\prod_v2^{-(2s_v-e_v)},\qquad
 \frac{w_{xy}(\omega)}{q_{xy}(\omega)}
      =N^{-k}2^{\sum_vC_v}.
\]
The last identity uses $\sum_vs_v=kd$. At the leaves the routing is determined. Summing unused switch coins contributes one, so $w_{xy}(\omega)$ is exactly the probability of the corresponding endpoint-and-coloring event. Thus
\[
 K_N(x,y)=\sum_\omega w_{xy}(\omega)
       =\frac1{(N)_k}\E_{q_{xy}}2^X.
\]
For the certificate alternative in Section~\ref{sec:certificate}, one may truncate the routing event itself. Define $A(x,y)=\sum_\omega w_{xy}(\omega)\mathbf1_{\{X\le L\}}$ and $E=K_N-A$. Then $A(x,y)\le2^L/(N)_k$, while
\begin{equation}\label{net:discarded-actual-law}
 \sum_y E(x,y)=\sum_{y,\omega}w_{xy}(\omega)\mathbf1_{\{X>L\}}
             =\PP_{\mathrm{actual},x}(X>L).
\end{equation}
Indeed $(y,\omega)$ partitions the actual switch experiment. The tail estimates below therefore bound discarded mass under the actual network law, without identifying that law with the auxiliary uniform-coloring law. The factor $(N)_k/N^k$ is essential: the endpoints form an injection, not a tuple sampled with replacement.

\begin{lemma}[Truncation and repeated irreducible copies]\label{net:truncation}
Let a finite group act transitively on a set of size $M$, and let $Q$ be the symmetric Markov kernel induced on its permutation module by a probability measure on that group. Suppose $Q=A+E$ entrywise, where $A,E\ge0$, every entry of $A$ is at most $B/M$, and the row and column sums of $E$ are at most $p$. If an irreducible representation occurs with multiplicity $m>0$, then
\[
 \|Q(\lambda)\|_{\op}\le \sqrt{B/m}+p.
\]
It is not necessary that $A$ preserve the irreducible subspaces.
\end{lemma}
\begin{proof}
Schur's test gives $\|E\|_{\op}\le p$. Since the row sums of $A$ are at most one,
$\|A\|_{\HS}^2\le(B/M)\sum_{x,y}A(x,y)\le B$.
Choose one maximizing unit vector for $Q(\lambda)$ in each of its $m$ orthogonal equivalent copies. They are orthonormal and each image under $A$ has norm at least $\|Q(\lambda)\|_{\op}-p$. Their squared image norms sum to at most $\|A\|_{\HS}^2$. This proves the claim.
\end{proof}
The first proof in Section~\ref{sec:factorial} instead truncates the endpoint density: it retains an entry when $(N)_kK_N(x,y)\le z$. Its density moment bounds the discarded row mass by Markov's inequality, and symmetry bounds the column mass. Both cutoffs use Lemma~\ref{net:truncation}; the letter $Q$ in that abstract lemma denotes any kernel satisfying its hypotheses.

For the routing-event cutoff above, reversal preserves the cutoff event, so a uniform discarded-row bound also bounds columns. Lemma~\ref{net:truncation} therefore gives $2^{L/2}/\sqrt m+p$.


\subsection{How often can selected paths meet?}
The next deterministic fact will control the information revealed about a closing route by other paths. Its logarithm is to base two.
\begin{lemma}[Butterfly encounters]\label{net:encounters}\label{n:butterfly-contacts}
Any $h$ distinct routed paths through a butterfly that updates each of $d$ binary coordinates once have at most $h\log_2h/2$ shared switches in total, with $0\log_2 0=0$. Their contact graph has maximum degree at most $d$. Counting both path endpoints at each shared switch gives at most $h\log_2h$ contacts.
\end{lemma}
\begin{proof}
Split by the coordinate edited last. If the two classes have sizes $a,h-a$, that layer contributes at most $\min(a,h-a)$ meetings. Apply induction to the earlier subnetworks and use
\[
 2\min(a,h-a)\le h\log_2h-a\log_2a-(h-a)\log_2(h-a).
\]
This is the binary entropy inequality $H_2(x)\ge2\min(x,1-x)$, where $H_2(x)=-x\log_2x-(1-x)\log_2(1-x)$. Empty terms are zero. Summing over the recursive splits proves the assertion. Distinct paths meet at most once: after a common switch they differ in a coordinate that is never edited again. A path meets at most one other path at each of its $d$ switches, proving the degree assertion.
\end{proof}

